\documentclass{article}
\usepackage{amsmath,amssymb}
\pagestyle{plain}
\begin{document}

Displays set their pieces in an order the source does not write, so this
document is mostly displays, each written with one row to a source line.

\begin{align}
  E &= \sum_{i=1}^{N} \alpha_i x_i + \frac{1}{2} \beta \\
  F &= \prod_{k=0}^{M} \gamma_k y_k - \frac{3}{4} \delta \\
  G &= \int_{0}^{T} \lambda(t) \, dt
\end{align}

The rows above are the energy, the force and the action.
\begin{gather}
  p = m v \\
  L = r \times p \\
  \tau = \frac{dL}{dt}
\end{gather}

A long one, broken by hand:
\begin{multline}
  H = \sum_{a} \frac{P_a^2}{2 M_a} + \sum_{i} \frac{p_i^2}{2 m} \\
  + \sum_{a<b} \frac{Z_a Z_b}{R_{ab}} - \sum_{a,i} \frac{Z_a}{r_{ai}} \\
  + \sum_{i<j} \frac{1}{r_{ij}}
\end{multline}

And a split inside one numbered equation:
\begin{equation}
  \begin{split}
    u &= w + z \\
      &= q \cdot s
  \end{split}
\end{equation}

A function by cases:
\begin{equation}
  f(x) =
  \begin{cases}
    x^2 & \text{if } x \ge 0 \\
    -x & \text{otherwise}
  \end{cases}
\end{equation}

Two short rows that share their letters:
\begin{align}
  a &= b + c \\
  c &= b + a
\end{align}

\end{document}
